% Moved here from the "Tests of the Standard Model" chapter: quark masses are a
% determination of Lagrangian parameters, not a test, and they use the same
% non-perturbative renormalisation machinery as the coupling.
\section{Quark masses}
\label{sec:quark-masses}

\todo{Section outline drafted mechanically --- prose and citations still to be
written. Delete this note once the section is yours.}

\todo{Write: what a quark mass \emph{is}. The bare lattice mass is not a
physical quantity; a renormalised mass exists only in a scheme and at a scale,
and the lattice computes the scheme-independent combinations --- the RGI mass,
and the ratios $m_s/m_{ud}$, $m_u/m_d$, $m_c/m_s$ --- from which any scheme
value follows. Emphasise that the ratios are cleaner than the absolute values
because the renormalisation constant cancels.}

\todo{Write: the tuning problem. Bare masses are fixed by matching a chosen set
of hadron masses, one per flavour, so the quark masses are inputs to that
procedure and outputs of nothing until the scale is set --- connect to the
scale-setting discussion of Ch.~\ref{ch:gauge-fields}, \S\ref{sec:gradient-flow},
and to the parameter-tuning walkthrough in Ch.~\ref{ch:data-analysis}.}

\todo{Write: renormalisation. Two routes, both non-perturbative and both
belonging to this chapter: the Schr\"odinger functional with step scaling,
which runs the mass over two orders of magnitude in scale before matching to
perturbation theory (the machinery of the ``Schr\"odinger-functional coupling''
section above, which wants a label), and RI-MOM
or RI-SMOM, which is cheaper but has a window problem. Mention gradient-flow
schemes as the newer alternative. The mass anomalous dimension and the
scheme conversion to $\overline{\rm MS}$.}

\todo{Write: the light quarks. Why $m_u$ is the hardest --- it requires either
isospin-breaking effects or QED, so it cannot be obtained from an isospin
symmetric calculation (cross-reference Ch.~\ref{ch:isospin-qed}), and why the
$m_u = 0$ solution to strong $CP$ is excluded by lattice determinations. This
last point is a genuine physics result and worth stating as one.}

\todo{Write: the heavy quarks. $m_c$ and $m_b$, the difficulty that $am_b > 1$
on accessible lattices, and the three responses --- effective theories (HQET,
NRQCD), relativistic actions at fine lattice spacing, and the step-scaling /
mass-extrapolation approaches. Also the current-current correlator method,
which extracts $m_c$ from moments and shares its machinery with the $\alpha_s$
determination earlier in this chapter.}

\todo{Write, briefly: what the masses are used for, as a pointer to
Ch.~\ref{ch:tests-of-sm} --- they enter the CKM extractions, the chiral
expansion, and every prediction that is compared to experiment. Refer the
reader to FLAG for current values rather than quoting numbers that will date.}
